% ============================================================
%  Chapter I · B 块：§I.5–§I.9（PDF p.23–p.36）
%  译注体例：BRIEF.md §4；本文件只写 \section，不写 \chapter
% ============================================================
% ---- 本文件局部宏（统一 fs 前缀，避免与其它块冲突）----
\providecommand{\fsSat}{\operatorname{Sat}}
\providecommand{\fsPerf}{\operatorname{Perf}}
\providecommand{\fsZ}{\mathrm{Z}}
\providecommand{\fsBZ}{D_{\mathrm{BZ}}}
\providecommand{\fslis}{D_{\mathrm{lis}}}
\providecommand{\fsrt}{\mathrm{rt}}
\providecommand{\fsdim}{\operatorname{dim}}
\providecommand{\fsHom}{\operatorname{Hom}}
\providecommand{\fsblk}{\blacksquare}
\providecommand{\fsCT}{\operatorname{CT}}
\providecommand{\fsHck}{\mathcal{H}\mathrm{ck}}
\providecommand{\Spd}{\operatorname{Spd}}
\providecommand{\fsL}{\mathbb{L}}

% ============================================================
% p.23 上半为 §I.4 尾部（A 块范围止于 p.22；若 A 块已译此段，删去本页第一段即可）
% ============================================================

\srcpageinl{23}每个这样的截面，我们得到 $S'$ 上的向量丛 $s^*T_{Z/X_S}$，其中 $T_{Z/X_S}$ 是切丛。粗略地说，$S'\to\cM_Z$ 的形变（即 $X_{S'}\to Z$ 在 $X_{S'}\to X_S$ 之上的形变）应当对应于整体截面 $H^0(X_{S'},s^*T_{Z/X_S})$，而其障碍应当属于 $H^1(X_{S'},s^*T_{Z/X_S})$。如果 $s^*T_{Z/X_S}$ 处处只有正的 Harder–Narasimhan 斜率，那么它在 $S'$ 上局部消失。与经典情形作类比，我们于是期待开子空间
\[
\cM_Z^{\mathrm{sm}}\subset\cM_Z,
\]
即 $s^*T_{Z/X_S}$ 具有正的 Harder–Narasimhan 斜率之处，在 $S$ 上（上同调）光滑。我们的关键几何结果证实了这一点，至少当 $Z\to X_S$ 拟投射时如此。

\begin{theorem}% 原编号 I.4.2
设 $Z\to X_S$ 局部地在 $S$ 上可嵌入（作为 adic 空间）$\PP^n_{X_S}$ 的某个开子集的 Zariski 闭子集。则 $\cM_Z\to S$ 在局部空间菱形中可表、可紧化，且 dim.\ trg 局部有限。此外，开子集 $\cM_Z^{\mathrm{sm}}\subset\cM_Z$ 在 $S$ 上上同调光滑。\yizhu{“dim.\ trg”是原文全书沿用的缩写（p.25、p.29 及 §IV 各处同此），此处按原样保留不译；它指 [Sch17a] 中引入的维数概念——原文 §IV 把它与 Krull 维数、上同调维数并列为三种维数——其精确定义见 [Sch17a]，本译本未作改动。}{A:p.25 “of finite dim.\ trg”、p.29 “dim.\ trg $f<\infty$” 同用此缩写；§IV 三种维数并列见 p.41}
\end{theorem}



在应用中，空间 $Z\to X_S$ 将是旗簇，它参数化 $X_S$ 上给定 $G$-主丛 $\mathcal{E}$ 的 $\QQ$-滤过。此时 $\cM_b$ 将是 $\cM_Z^{\mathrm{sm}}$ 的一个开子集。

定理 I.4.2 的证明需要若干项创新。第一项是形式光滑性的概念，其中（在完满 setting 中并不存在的）无穷小加厚被小 étale 邻域取代。由此得到的概念与经典拓扑中绝对邻域收缩核的概念 [Bor67] 有密切关系。我们证明，几乎所有上同调光滑映射的例子也是形式光滑的，包括 Banach–Colmez 空间与 $\Bun_G$。我们还证明 $\cM_Z^{\mathrm{sm}}\to S$ 是形式光滑的，这归结为一些精细的估计：对 affinoid 完满空间的任意 Zariski 闭浸入 $T_0\subset T$，把截面 $X_{T_0}\to Z$ 展布到 $T_0\subset T$ 的小邻域上——这里我们关键地用到“所有 Harder–Narasimhan 斜率均为正”的假设。再结合“Zariski 闭蕴含强 Zariski 闭”这一定理 [BS22]，就使我们能在相差（上同调与形式）光滑映射的意义下，把 $\cM_Z^{\mathrm{sm}}$ 写成某个在 $S$ 上的球上 étale 的空间的收缩。在经典拓扑中，这当然不足以保证上同调光滑——坐标十字是 $\RR^2$ 的收缩——但它确实蕴含常层 $\FF_\ell$ 在 $S$ 上普遍局部无环。出于这个原因，以及为了层在 $\Bun_G$ 上以及几何 Satake 的其它应用，我们于是也在自己的 setting 中发展一套普遍局部无环层的一般理论。为完成证明，我们用一个形变到法锥的论证来说明，对偶化复形与 Banach–Colmez 空间 $\mathrm{BC}(s^*T_{Z/X_S})$ 的对偶化复形“相同”。

% ============================================================
\section{$\Bun_G$ 上的 $\ell$-进层}
% ============================================================

为得到我们的结果，需要定义 $\Bun_G$ 上的 $\ell$-进层范畴。更精确地说，我们将对每个 $\ZZ_\ell$-代数 $\Lambda$ 定义一个范畴
\[
D(\Bun_G,\Lambda),
\]
即 $\Bun_G$ 上 $\Lambda$-模层的范畴。若 $\Lambda$ 被 $\ell$ 的某个幂零化，这样的定义是 [Sch17a] 的主要成就。但我们主要关心的是 $\Lambda=\ol{\QQ}_\ell$ 的情形。在概形（代数闭域上有限型）的情形，从挠系数过渡到 $\ol{\QQ}_\ell$-系数\srcpageinl{24} 大体上是形式的：粗略地说，
\[
D(X,\ol{\QQ}_\ell)=\operatorname{Ind}\Bigl(\varprojlim_n D^b_c(X,\ZZ/\ell^n\ZZ)\otimes_{\ZZ_\ell}\ol{\QQ}_\ell\Bigr).
\]
然而这一约定背后是关于构造层的强有限性结果；特别地，构造层之间的态射空间是有限的。对 $\Bun_G$，或对光滑表示范畴而言，仍然存在紧对象（在光滑表示的情形下由紧诱导表示给出），但它们的自同态代数是 Hecke 代数，而后者是无限维的。沿同样的思路作定义，就会把所有 Hecke 代数换成它们的 $\ell$-进完备化，这将剧烈改变表示范畴。

我们对 $D(\Bun_G,\Lambda)$ 的定义则一般地涉及一些新想法，用到 Clausen–Scholze [CS] 在 pro-étale（或 v-）site 情形下发展的固体模思想\yizhu{“固体模”（solid module）是 Clausen–Scholze 凝聚数学中的概念：凝聚模里满足某种完备性条件者称为固体的；本译本沿用裁定，不译 solid。}{B:Clausen–Scholze《Condensed Mathematics》与 [CS]}；最终，$D(\Bun_G,\Lambda)$ 被定义为某个满子范畴
\[
D_{\mathrm{lis}}(\Bun_G,\Lambda)\subset D_{\fsblk}(\Bun_G,\Lambda),
\]
即 $\Bun_G$ 的 v-site 上 $\Lambda$-模的固体复形范畴 $D_{\fsblk}(\Bun_G,\Lambda)$ 的某个满子范畴。固体层的形式体系（其想法归功于 Clausen 与第二作者）在第七章中发展。它带来一些有趣的意外；特别地，拉回 $f^*$ 总有一个左伴随 $f_\natural$，后者满足基变换与一个投射公式。（作为回报，$Rf_!$ 一般不再存在。）


\begin{theorem}% 原编号 I.5.1
设 $\Lambda$ 为任意 $\ZZ_\ell$-代数。
\begin{enumerate}
\item[(i)] 经由切除三角形，存在 $D(\Bun_G,\Lambda)$ 按各个 $D(\Bun^b_G,\Lambda)$（\\mbox{$b\in B(G)$}）的无穷半正交分解。
\item[(ii)] 对每个 \\mbox{$b\in B(G)$}，沿
\[
\Bun^b_G\cong[*/\wt{G}_b]\longrightarrow[*/G_b(E)]
\]
的拉回给出一个等价
\[
D([*/G_b(E)],\Lambda)\cong D(\Bun^b_G,\Lambda),
\]
而 $D([*/G_b(E)],\Lambda)\cong D(G_b(E),\Lambda)$ 等价于 $G_b(E)$ 在 $\Lambda$-模上的光滑表示范畴的导出范畴。
\item[(iii)] 范畴 $D(\Bun_G,\Lambda)$ 是紧生成的，且复形 $A\in D(\Bun_G,\Lambda)$ 是紧的，当且仅当对所有 \\mbox{$b\in B(G)$}，限制
\[
i^{b*}A\in D(\Bun^b_G,\Lambda)\cong D(G_b(E),\Lambda)
\]
是紧的，并且对几乎所有 $b$ 为零。这里，$D(G_b(E),\Lambda)$ 中的紧性等价于落在由 $\operatorname{c-Ind}_K^{G_b(E)}\Lambda$（$K$ 取遍 $G_b(E)$ 的开 pro-$p$-子群）生成的厚三角子范畴中。
\item[(iv)] 在紧对象的子范畴 $D(\Bun_G,\Lambda)^\omega\subset D(\Bun_G,\Lambda)$ 上，存在一个 Bernstein–Zelevinsky 对偶函子
\[
\fsBZ:(D(\Bun_G,\Lambda)^\omega)^{\mathrm{op}}\longrightarrow D(\Bun_G,\Lambda)^\omega,
\]
并带有函子性的等同
\[
R\Hom(A,B)\cong\pi_\natural\bigl(\fsBZ(A)\otimes^{\fsL}_\Lambda B\bigr).
\]
\srcpageinl{25}其中 $B\in D(\Bun_G,\Lambda)$，而 $\pi:\Bun_G\to *$ 为投影。函子 $\fsBZ$ 是一个等价，且 $\fsBZ^2$ 自然同构于恒等函子。它与 $D(G_b(E),\Lambda)$（\\mbox{$b\in B(G)$} 为 basic 元）上通常的 Bernstein–Zelevinsky 对偶相容。
\item[(v)] 对象 $A\in D(\Bun_G,\Lambda)$ 是普遍局部无环的（相对于 $\Bun_G\to *$），当且仅当对所有 \\mbox{$b\in B(G)$}，限制
\[
i^{b*}A\in D(\Bun^b_G,\Lambda)\cong D(G_b(E),\Lambda)
\]
是容许的，即对所有开 pro-$p$-子群 $K\subset G_b(E)$，复形 $(i^{b*}A)^K$ 是完美的。普遍局部无环复形在 Verdier 对偶下保持不变，并满足 Verdier 双对偶。
\end{enumerate}
\end{theorem}


这一定理把表示论中的许多基本概念——有限表现对象、容许表示、Bernstein–Zelevinsky 对偶、光滑对偶——推广到 $D(\Bun_G,\Lambda)$ 的 setting 中。

当 $\Lambda$ 是 $\ell$-幂挠时，部分 (i) 与 (ii) 是容易的。一般情形下，它们的证明涉及 $D(\Bun_G,\Lambda)=D_{\mathrm{lis}}(\Bun_G,\Lambda)$ 的精确定义，并且有些微妙。部分 (iii) 用到 $i^{b*}$ 有左伴随，后者于是自动保持紧对象（从而给出紧生成元）。利用图表
\[
\begin{tikzcd}
{[*/G_b(E)]} & \cM_b \arrow[l,"q_b"'] \arrow[r,"\pi_b"] & \Bun_G,
\end{tikzcd}
\]
这个左伴随定义为 $\pi_{b\natural}q_b^*$。验证它确实是左伴随，在某种意义上归结为如下断言：$\cM_b$ 沿闭子空间 $[*/G_b(E)]\subset\cM_b$ 是“严格局部”的，即对所有 $A\in D(\cM_b,\Lambda)$，限制映射
\[
R\Gamma(\cM_b,A)\longrightarrow R\Gamma([*/G_b(E)],A)
\]
是同构。这建立在关于 $\cM_b$ 的拓扑性质的细致分析之上，特别是 $\wt{\cM}_b\setminus\{*\}$ 是空间菱形这一事实，以及下面的定理 I.5.2。至于部分 (iv)，(iii) 中的构造在一类生成元上给出了 $\fsBZ(A)$ 的存在性，从而在一般情况下也给出；而与 (iii) 中类似的论证证明了双对偶性。最后，部分 (v) 本质上是形式推论。\yizhu{“（当 $\Lambda$ 是 $\ell$-幂挠时）容易”与“本质上是形式推论”两处原文均未写出论证。前者：$\Lambda$ 为 $\ell$-幂挠时可用 [Sch17a] 的六函子形式体系，闭浸入的推出与拉回给出切除三角形，从而逐步得到半正交分解；后者：由 (iii) 中 $i^{b*}$ 的左伴随保持紧性，配合 (v) 中“容许”的定义即可逐项验证。}{A:p.25 “are easy when $\Lambda$ is $\ell$-power torsion”/“is essentially a formal consequence”}



证明所需的关键上同调结果如下，它施用于 $\wt{\cM}_b\setminus\{*\}$（或它的商）。该结果玩味的是点 $*$ 并非拟分离这一微妙之处。

\begin{theorem}% 原编号 I.5.2
设 $X$ 为空间菱形，使得 $f:X\to *$ 部分紧合，且 dim.\ trg 有限。则对任意 affinoid 完满空间 $S$，基变换 $X_S=X\times S$ 自然地具有两个端。取关于其中一个端的紧支上同调（而在另一端不加任何支撑条件），则对一切 $A\in D^+_{\fsblk}(X,\ZZ_\ell)$（相应地，若 $f$ 是 $\ell$-上同调光滑的，则对一切 $A\in D_{\fsblk}(X,\ZZ_\ell)$）有
\[
R\Gamma_{\partial\text{-}c}(X_S,A)=0.
\]
\end{theorem}

举例来说，若 $X=\Spa\ol{\FF}_q((t))$，则 $X_S=\mathbb{D}^*_S$ 是 $S$ 上的开单位圆盘，其两个端分别是原点与边界，于是有
\[
R\Gamma_{\partial\text{-}c}(\mathbb{D}^*_S,\ZZ_\ell)=0.
\]
特别地，$\Spa\ol{\FF}_q[[t]]$ 的上同调与闭点上的截面一致，这表明 $\Spa\ol{\FF}_q[[t]]$ 是“严格局部”的。同样的现象也在 $\cM_b$ 处发挥作用。

% ============================================================
\section{几何 Satake 等价}
% ============================================================

\srcpageinl{26}为了定义 Hecke 算子，我们需要证明几何 Satake 等价，即把对偶群 $\hat G$ 的表示对应为局部 Hecke 叠上的层。而为了分析 Hecke 算子的复合，事实上还必须分析多个点处的修正。

于是，对任意有限集 $I$，我们考虑参数化次数为 $1$ 的 Cartier 除子 $D_i\subset X_S$（$i\in I$）的模空间 $(\Div^1)^I$。局部地在 $S$ 上，每个 $D_i$ 定义了 $S$ 在 $E$ 上的一个 untilt $S_i^\sharp$，并且可以构造 $\cO_{X_S}$ 沿诸 $D_i$ 之并的完备化 $B^+$。把 $D_i$ 取逆就定义了 $B^+$ 的一个局部化 $B$。于是可以定义正环路群 $L^+_{(\Div^1)^I}G$ 与环路群 $L_{(\Div^1)^I}G$，其取值分别由 $G(B^+)$ 与 $G(B)$ 给出；为简洁起见，我们在这里就简单地写作 $L^+G$ 与 $LG$。然后可以定义局部 Hecke 叠
\[
\fsHck^I_G=[L^+G\backslash LG/L^+G]\longrightarrow(\Div^1)^I.
\]
对 $d=|I|$，它事实上已经定义在次数为 $d$ 的 Cartier 除子的模空间 $\Div^d=(\Div^1)^d/\Sigma_d$ 上。我们常常打破对称性，先对右边取商，以定义 Beilinson–Drinfeld Grassmann 簇
\[
\Gr^I_G=LG/L^+G\longrightarrow(\Div^1)^I,
\]
从而 $\fsHck^I_G=L^+G\backslash\Gr^I_G$。

Beilinson–Drinfeld Grassmann 簇 $\Gr^I_G\to(\Div^1)^I$ 是一个小 v-层，它可以写成递增的闭子层之并，而这些闭子层通过界定相对位置而是紧合的、且在空间菱形中可表；这是 [SW20] 的主要结果之一。另一方面，$L^+G$ 可以写成截断正环路群的逆向极限，这些截断正环路群在局部空间菱形中可表并且上同调光滑；此外，在每个有界子集上，它通过这样一个有限维商作用。这本质上把 $\fsHck^I_G$ 的所有有界子集的研究归约到 Artin 叠。

特别地，通过界定相对位置，可以把局部 Hecke 叠写成在 $(\Div^1)^I$ 上拟紧的闭子叠的递增并。以下我们假定系数 $\Lambda$ 被 $\ell$ 的某个幂零化，以便使用 [Sch17a] 的理论。令
\[
D_{\text{\'et}}(\fsHck^I_G,\Lambda)^{\mathrm{bd}}\subset D_{\text{\'et}}(\fsHck^I_G,\Lambda)
\]
为在 $(\Div^1)^I$ 上具有拟紧支撑的所有对象构成的满子范畴。它在卷积 $\star$ 下是幺半范畴。这里我们使用卷积图表
\[
\begin{tikzcd}
\fsHck^I_G\times_{(\Div^1)^I}\fsHck^I_G & L^+G\backslash LG\times_{L^+G}LG/L^+G \arrow[l,"{(p_1,p_2)}"'] \arrow[r,"m"] & \fsHck^I_G
\end{tikzcd}
\]
并定义
\[
A\star B=Rm_*\bigl(p_1^*A\otimes^{\fsL}_{\Lambda}p_2^*B\bigr).
\]
映射 $m$ 是 ind-紧合的（其纤维是 $\Gr^I_G$），从而在任何有界子集上紧合；因此紧合基变换保证这定义了一个结合的幺半结构。

在 $D_{\text{\'et}}(\fsHck^I_G,\Lambda)^{\mathrm{bd}}$ 上，可以定义一个相对反常 $t$-结构（其中对象反常，当且仅当它在 $(\Div^1)^I$ 的每个几何纤维上都反常）。对这个 $t$-结构，卷积\srcpageinl{27} $\star$ 是左 $t$-恰当的（$t$-恰当性只在涉及 $\Lambda$ 上非平坦性的问题上失效）。为证明 $t$-结构良定义以及反常性在卷积下保持，我们移植了 Braden 关于双曲局部化的定理 [Bra03]，以及一个到 Witt 向量仿射 Grassmann 簇的退化 [Zhu17]、[BS17]。我们将在下文进一步讨论双曲局部化。

我们指出，$p$-进几何中并没有反常层的一般理论，问题在于给 adic 空间的一个点无歧义地指定维数是困难的（关于点的拓扑超越度已知的结果，以及特征 $p$ 情形下特有的微妙之处，可参看 [Tem21]）。特别地，我们不知道如何在一般情况下于 $(\Div^1)^I$ 上定义反常层的概念，这正是我们退而只在纤维中要求反常性的原因。这里我们用到：叠 $\fsHck^I_G\to(\Div^1)^I$ 的所有几何纤维都只有可数多个点，且这些点可以用支配余特征标 $\mu_i$ 显式列举，而相应开 Schubert 胞的维数 $\sum_i\langle 2\rho,\mu_i\rangle$ 可以手工指定。

\begin{remark}% 原编号 I.6.1
受此启发，我们意识到：对概形之间任意局部有限型的映射 $f:X\to S$，都可以定义一个相对反常 $t$-结构，其相对反常性等价于在所有几何纤维上的反常性，参看 [HS23]。
\end{remark}

此外，还可以限制到那些在 $(\Div^1)^I$ 上普遍局部无环的复形 $A\in D_{\text{\'et}}(\fsHck^I_G,\Lambda)^{\mathrm{bd}}$。这一条件同样在卷积下保持。

\begin{definition}% 原编号 I.6.2
\emph{Satake 范畴}
\[
\fsSat^I_G(\Lambda)\subset D_{\text{\'et}}(\fsHck^I_G,\Lambda)^{\mathrm{bd}}
\]
是所有满足以下条件的 $A\in D_{\text{\'et}}(\fsHck^I_G,\Lambda)^{\mathrm{bd}}$ 构成的范畴：$A$ 反常；$A$ 在 $\Lambda$ 上平坦（即对所有 $\Lambda$-模 $M$，$A\otimes^{\fsL}_{\Lambda}M$ 也反常）；且 $A$ 在 $(\Div^1)^I$ 上普遍局部无环。\yizhu{“在 $\Lambda$ 上平坦”在此按层论意义理解：张量积 $A\otimes^{\fsL}_\Lambda M$ 的反常性对一切 $\Lambda$-模 $M$ 成立，等价于 $A$ 的每一阶上同调层在 $\Lambda$ 上平坦；原文括注已给出这一定义，此处补出等价形式以便与下文 ULA 条件对照。}{I:推理，配合定义 I.6.2 括注}
\end{definition}



直观上，$\fsSat^I_G(\Lambda)$ 是“$\fsHck^I_G\to(\Div^1)^I$ 上反常层的平坦族”，其中平坦既指在 $(\Div^1)^I$ 上的几何意义的平坦（由普遍局部无环性编码），也指系数 $\Lambda$ 中的代数意义的平坦。Satake 范畴 $\fsSat^I_G(\Lambda)$ 在卷积下是幺半范畴。此外，它对 $I$ 是共变函子性的。

事实上，这一幺半结构自然升级为对称幺半结构。这依赖于融合积，而对此允许一般的有限集 $I$ 至关重要。具体地说，给定有限集 $I_1,\dots,I_k$，令 $I=I_1\sqcup\dots\sqcup I_k$，则有同构
\[
\fsHck^I_G\times_{(\Div^1)^I}(\Div^1)^{I;I_1,\dots,I_k}\cong\prod_{j=1}^{k}\Bigl(\fsHck^{I_j}_G\times_{(\Div^1)^{I_j}}(\Div^1)^{I;I_1,\dots,I_k}\Bigr),
\]
其中 $(\Div^1)^{I;I_1,\dots,I_k}\subset(\Div^1)^I$ 是使下述条件成立的开子集：只要 $i,i'\in I$ 落在不同的 $I_j$ 中，就有 $x_i\neq x_{i'}$。外张量积于是定义一个函子
\[
\boxtimes_{j=1}^{k}:\prod_{j=1}^{k}\fsSat^{I_j}_G(\Lambda)\longrightarrow\fsSat^{I;I_1,\dots,I_k}_G(\Lambda),
\]
其中 $\fsSat^{I;I_1,\dots,I_k}_G(\Lambda)$ 是 $\fsSat^I_G(\Lambda)$ 对于 $\fsHck^I_G\times_{(\Div^1)^I}(\Div^1)^{I;I_1,\dots,I_k}$ 的变体。然而，限制函子
\[
\fsSat^I_G(\Lambda)\longrightarrow\fsSat^{I;I_1,\dots,I_k}_G(\Lambda)
\]
\srcpageinl{28}是全忠实的，且外张量积的本质像落在其本质像之中。于是我们得到自然的函子
\[
\ast_{j=1}^{k}:\prod_{j=1}^{k}\fsSat^{I_j}_G(\Lambda)\longrightarrow\fsSat^{I}_G(\Lambda),
\]
它与诸 $I_j$ 的排序无关。特别地，对任意 $I$ 我们得到函子
\[
\fsSat^I_G(\Lambda)\times\fsSat^I_G(\Lambda)\longrightarrow\fsSat^{I\sqcup I}_G(\Lambda)\longrightarrow\fsSat^{I}_G(\Lambda),
\]
这里用到 $\fsSat^I_G(\Lambda)$ 对 $I$ 的函子性；它定义了 $\fsSat^I_G(\Lambda)$ 上一个与 $\star$ 交换的对称幺半结构 $\ast$。这称为融合积。一般地，对任何带有交换的幺半结构 $\star$ 的对称幺半范畴 $(\mathcal{C},\ast)$，幺半结构 $\ast$ 必然与 $\star$ 一致；因此融合积细化卷积积。（正如几何 Satake 中通常的做法，我们实际上需要略微改动 $\ast$：在交换性约束中引入某些符号，符号取决于反常层支撑的奇偶性。）\yizhu{“必然与 $\star$ 一致”：两个幺半结构互相交换时，其中一个的交换性约束给出二者间的自然同构，从而 $(C,\ast)$ 的幺半积与 $\star$ 在同构意义下相同；原文未给论证，此为标准 Eckmann–Hilton 型论证。}{I:推理（Eckmann–Hilton 论证）}



此外，把 $A\in\fsSat^I_G(\Lambda)$ 限制到 $\Gr^I_G$ 再前推到 $(\Div^1)^I$，则所有上同调层都是 $(\Div^1)^I$ 上 $\Lambda$-模的局部系统。由 Drinfeld 引理的一个版本，它们等价于 $W_E^I$ 在 $\Lambda$-模上的表示。这定义了一个对称幺半纤维函子
\[
F^I:\fsSat^I_G(\Lambda)\longrightarrow\Rep_{W_E^I}(\Lambda),
\]
其中 $\Rep_{W_E^I}(\Lambda)$ 是 $W_E^I$ 在有限投射 $\Lambda$-模上的连续表示范畴。利用 Tannaka 对偶的一个版本，就可以在 $\Rep_{W_E^I}(\Lambda)$ 的 Ind-范畴中构造一个 Hopf 代数，使得 $\fsSat^I_G(\Lambda)$ 由它的表示范畴给出（在 $\Rep_{W_E^I}(\Lambda)$ 内部）。对任意有限集 $I$，它由 $I=\{*\}$ 情形对应的 Hopf 代数的 $I$ 个副本的张量积给出；而后者由一个带 $W_E$-作用的、$\Lambda$ 上的仿射群概形 $\wt{G}$ 给出。

\begin{theorem}% 原编号 I.6.3
存在与朗兰兹对偶群的典范同构 $\wt{G}\cong\hat G_\Lambda$，在此同构下，$W_E$ 在 $\wt{G}$ 上的作用与 $W_E$ 在 $\hat G$ 上的通常作用相差一个显式的分圆扭转。若 $\sqrt q\in\Lambda$，则分圆扭转可以被平凡化，而 $\fsSat^I_G(\Lambda)$ 自然等价于 $(\hat G\rtimes W_E)^I$ 在有限投射 $\Lambda$-模上的表示范畴。
\end{theorem}

这一定理因而是 Mirković–Vilonen 的定理 [MV07] 的一个版本，并配上 Gaitsgory [Gai07] 对一般 $I$ 的加细。（我们指出，我们叙述的是对任意 $\Lambda$（不必正则）都成立的定理；这样的表述似乎不见于文献。此外，我们给出的是纯局部的证明：大多数证明需要在（通常的）曲线上作整体化。）与 Mirković–Vilonen 相反，我们实际上构造了 $\wt{G}$ 的一个显式钉住。至于证明，可以先将 $\Lambda$ 限制为 $\ZZ/\ell^n\ZZ$；对 $n$ 取极限，就能真正构造出 $\ZZ_\ell$ 上的群概形。它的一般纤维是约化的，因为 $\QQ_\ell$-系数的 Satake 范畴是（几何）半单的：为此我们再次使用到 Witt 向量仿射 Grassmann 簇的退化，以及概形的分解定理。为识别这个约化群，我们先对环面论证，再对秩 $1$ 群论证，此时一切都归约到 $G=\mathrm{PGL}_2$，而后者借极小 Schubert 胞容易分析。此处的钉住包含一个分圆扭转，因为 $\Gr_{\mathrm{PGL}_2}$ 的极小 Schubert 簇 $\PP^1$ 的上同调当然含有分圆扭转。此后，我们应用双曲局部化，对 $G$ 的任意 Levi 子群 $M$ 构造对称幺半函子 $\fsSat_G\to\fsSat_M$，对偶地诱导出映射 $\wt{M}\to\wt{G}$。\srcpageinl{29}由此容易得到与 $\hat G_{\QQ_\ell}$ 的同构，并包含一个钉住。由于这些映射 $\wt{M}\to\wt{G}$ 甚至整地有定义，且由 Bruhat–Tits 理论 $\hat G(\ZZ_\ell)\subset\hat G(\QQ_\ell)$ 是由秩 $1$ Levi 子群生成的极大紧开子群，于是可以推出 $\wt{G}\cong\hat G$ 整地成立，同样带有一个显式的（分圆）钉住。

我们还需要下面这条关于一个自然对合的补记。即：局部 Hecke 叠 $\fsHck_G$ 有一个自然对合 $\mathrm{sw}$，由交换两个 $G$-主丛的角色给出；在用以 $LG$ 表述的展示中，它由 $LG$ 上的取逆诱导。于是 $\mathrm{sw}^*$ 自然地诱导出 $\fsSat_G(\Lambda)$ 的一个对合，而这个对合可以升级为与纤维函子交换的对称幺半函子，从而实现 $\wt{G}\cong\hat G$ 的一个 $W_E$-等变自同构。

\begin{proposition}% 原编号 I.6.4
$\mathrm{sw}^*$ 在 $\fsSat_G$ 上的作用诱导 $\hat G$ 的自同构，它是分裂群 $\hat G$ 的 Chevalley 对合经 $\hat\rho(-1)$ 共轭之后所得的自同构。
\end{proposition}

对我们全部论证都至关重要的是双曲局部化函子。在 Beilinson–Drinfeld Grassmann 簇的 setting 中，设 $P^+,P^-\subset G$ 是两个互为相反抛物子群，具有共同的 Levi 子群 $M$。我们得到一个图表
\[
\begin{tikzcd}
& \Gr^I_{P^+} \arrow[dl,"q^+"] \arrow[dr,"p^+"] & \\
\Gr^I_G & & \Gr^I_M \\
& \Gr^I_{P^-} \arrow[ul,"q^-"] \arrow[ur,"p^-"] &
\end{tikzcd}
\]
我们得到两个“常数项”函子
\[
\fsCT^+=R(p^+)_!(q^+)^*,\qquad
\fsCT^-=R(p^-)_*R(q^-)^!:\quad
D_{\text{\'et}}(\Gr^I_G,\Lambda)^{\mathrm{bd}}\longrightarrow D_{\text{\'et}}(\Gr^I_M,\Lambda)^{\mathrm{bd}},
\]
并且可以构造一个自然变换 $\fsCT^-\to\fsCT^+$。函子 $\fsCT^+$ 在经典情形下对应于 Satake 变换，即沿 $U^+$ 的幂幺根中的轨道作积分。双曲局部化断言：限制到 $L^+G$-等变对象上时，变换 $\fsCT^-\to\fsCT^+$ 是等价。这有许多推论；注意 $\fsCT^+$ 由左伴随函子构成，而 $\fsCT^-$ 由右伴随函子构成，因此若二者同构，双曲局部化便兼具二者之长。特别地，双曲局部化与所有余极限、所有极限交换，保持（相对）反常性、普遍局部无环性，与任意基变换交换，等等。\yizhu{原文此句末印作 “etc.\ .”，多出一个句点，应为排版笔误，译文按“等等”处理。}{A:p.29 末行 “commutes with any base change, etc. .”}



事实上，这只是下述更一般断言的一个特殊情形。设 $S$ 为任意小 v-叠，$f:X\to S$ 为紧合的、在空间菱形中可表且 $\fsdim$.\ trg $f<\infty$ 的映射。假设在 $X/S$ 上有一个 $\GG_m$-作用，其中 $\GG_m(R,R^+)=R^\times$。$\GG_m$-作用的不动点 $X^0\subset X$ 构成闭子叠。我们假设可以定义吸引轨迹 $X^+\subset X$ 与排斥轨迹 $X^-\subset X$，它们由局部闭子空间的无交并给出，并且在其上，\srcpageinl{30}$\GG_m$-作用在 $t\to0$（相应地 $t\to\infty$）时具有极限。我们得到一个图表
\[
\begin{tikzcd}
& X^+ \arrow[dl,"q^+"] \arrow[dr,"p^+"] & \\
X & & X^0, \\
& X^- \arrow[ul,"q^-"] \arrow[ur,"p^-"] &
\end{tikzcd}
\]
只要选取一个余特征标 $\mu:\GG_m\to G$，使其动力抛物子群为 $P^+,P^-$，此图就推广了上面的图表（的有界部分）。可以定义
\[
L^+=R(p^+)_!(q^+)^*,\qquad L^-=R(p^-)_*R(q^-)^!:\quad D_{\text{\'et}}(X,\Lambda)\longrightarrow D_{\text{\'et}}(X^0,\Lambda)
\]
以及一个自然变换 $L^-\to L^+$。下面是我们版本的 Braden 定理 [Bra03]，另可参看 [Ric19]。

\begin{theorem}% 原编号 I.6.5
限制到 $D_{\text{\'et}}(X/\GG_m,\Lambda)\to D_{\text{\'et}}(X,\Lambda)$ 的本质像上时，变换 $L^-\to L^+$ 是等价。\yizhu{“本质像上”一句即：设 $\pi:X\to X/\GG_m$，则对 $\pi$ 的像中的对象，$L^-$ 与 $L^+$ 同构；这正是双曲局部化在 $\Gr$ 情形（定理前的讨论）的一般化，故原文未重复论证。}{A:p.29 末段“双曲局部化断言……”}
\end{theorem}



证明用到下述原则：若 $Y\to S$ 部分紧合且带有 $\GG_m$-作用，使得商叠 $Y/\GG_m$ 在 $S$ 上 qcqs，则 $Y$ 同样具有两个端，并且对任意 $A\in D_{\text{\'et}}(Y/\GG_m,\Lambda)$，$Y$ 上系数在 $A$ 中的部分紧支撑上同调恒为零。

% ============================================================
\section{shtuka 模空间的上同调}
% ============================================================

至此我们已经定义了 $D(\Bun_G,\Lambda)$；利用几何 Satake 等价与图表
\[
\begin{tikzcd}
& \Hck^I_G \arrow[r,"q"] \arrow[dl,"h_1"'] \arrow[dr,"h_2"] & \fsHck^I_G \\
\Bun_G & & \Bun_G\times(\Div^1)^I
\end{tikzcd}
\]
就可以定义 Hecke 算子
\[
T_V=R h_{2*}\bigl(h_1^*(\,\cdot\,)\otimes^{\fsL}_\Lambda q^*\mathcal{S}_V\bigr):\quad
D(\Bun_G,\Lambda)\longrightarrow D\bigl(\Bun_G\times(\Div^1)^I,\Lambda\bigr)
\]
其中 $V\in\fsSat^I_G(\Lambda)$，而 $\mathcal{S}_V$ 是 $\fsHck^I_G$ 上相应的层。至少在 $\Lambda$ 被 $\ell$ 的某个幂零化时这行得通。事实上我们可以把这个函子延拓到任意 $\ZZ_\ell$-代数 $\Lambda$。此外，它的像落在那些沿 $(\Div^1)^I$ 方向局部常值的对象所构成的满子范畴中，从而给出一个函子
\[
T_V:D(\Bun_G,\Lambda)\longrightarrow D(\Bun_G,\Lambda)^{BW_E^I}
\]
\srcpage{31}取值于 $D(\Bun_G,\Lambda)$ 中 $W_E^I$-等变对象所成的范畴。其证明出人意料地形式化：先用一个归纳论证归约到 $I=\{*\}$，再利用 $\Div^1=\Spd\dot E/\varphi^\ZZ$ 仍然只是一个点。更精确地说，用到
\[
D(\Bun_G,\Lambda)\longrightarrow D(\Bun_G\times\Spd\dot E,\Lambda)
\]
是等价。\yizhu{原文的归纳论证没有展开，可补出的步骤是：对 $|I|$ 归纳，用融合积把 $I=I'\sqcup\{*\}$ 时的算子分解为两个因子的外积，$W_E^I=W_E^{I'}\times W_E$ 的作用也逐因子给出；单因子情形 $I=\{*\}$ 正是原文紧接着说明的那一步——$\Div^1=\Spd\dot E/\varphi^\ZZ$“仍然只是一个点”，即 $\Bun_G\times\Spd\dot E\to\Bun_G$ 为同构，故 $D(\Bun_G,\Lambda)\to D(\Bun_G\times\Spd\dot E,\Lambda)$ 是范畴等价，算子沿 $\Spd\dot E$ 方向常值，$\varphi$ 的作用给出 $\Div^1$ 方向的局部系统。}{I:推理（补原文省略的归纳步骤与“一个点”的含义）}



\begin{remark}% 原编号 I.7.1
为定义 $D(\Bun_G,\Lambda)^{BW_E^I}$，需要把 $D(\Bun_G,\Lambda)$ 升级为一个凝聚 $\infty$-范畴；此时它才是凝聚群 $W_E^I$ 的等变对象这一概念。
\end{remark}

我们结果的一个直接推论是：忘记 $W_E^I$-等变性之后，$T_V$ 保持有限性。注意由于几何 Satake 等价是幺半的，有 $T_V\circ T_W\cong T_{V\otimes W}$。这形式地蕴含 $T_V$ 是 $T_{V^*}$ 的左伴随兼右伴随。由此不难证明下述结果。

\begin{theorem}% 原编号 I.7.2
函子 $T_V:D(\Bun_G,\Lambda)\to D(\Bun_G,\Lambda)$ 保持紧对象与普遍局部无环对象。此外，它与 Bernstein–Zelevinsky 对偶和 Verdier 对偶交换，即存在自然同构
\[
\fsBZ(T_V(A))\cong T_{\mathrm{sw}^*V^\vee}(\fsBZ(A))
\quad\text{与}\quad
R\mathcal{H}om(T_V(A),\Lambda)\cong T_{(\mathrm{sw}^*V)^\vee}R\mathcal{H}om(A,\Lambda).
\]
\end{theorem}

这里 $\mathrm{sw}^*$ 是 $\fsSat^I_G$ 的对合，由命题 I.6.4，它由 $\hat G$ 的 Chevalley 对合经 $\hat\rho(-1)$ 共轭而诱导。

这一定理对 shtuka 模空间的上同调有具体推论。为简单起见，我们在此用被 $\ell^n$（对某个 $n$）零化的 $\Lambda$-代数中的系数来叙述；至于一般表述，我们需要更精确地讨论围绕导出范畴的基础性问题。在 [SW20, Lecture XXIII] 中，对任意一族余特征标的共轭类 $\{\mu_i\}_i$（其定义域为 $E_i/E$）以及 $b\in B(G)$，定义了局部 shtuka 模空间的塔
\[
f_K:(\Sht_{(G,b,\mu_\bullet),K})_{K\subset G(E)}\longrightarrow\prod_{i\in I}\Spd\dot E_i,
\]
其中 $K$ 取遍 $G(E)$ 的紧开子群；这些模空间带有相容的 étale 周期映射
\[
\pi_K:\Sht_{(G,b,\mu_\bullet),K}\longrightarrow\Gr^{\mathrm{tw}}_{G,\prod_{i\in I}\Spd\dot E_i,\le\mu_\bullet}.
\]
这里 $\Gr^{\mathrm{tw}}_{G,\prod_{i\in I}\Spd\dot E_i}\to\prod_{i\in I}\Spd\dot E_i$ 是卷积仿射 Grassmann 簇的某个扭转形式，参看 [SW20, Section 23.5]。令 $W$ 为最高权表示的外张量积 $\boxtimes_{i\in I}V_{\mu_i}$，$S_W$ 为 $\Gr^{\mathrm{tw}}_{G,\prod_{i\in I}\Spd\dot E_i}$ 上相应的层。我们仍以 $S_W$ 记它到 $\Sht_{(G,b,\mu_\bullet),K}$ 的拉回。

\begin{corollary}% 原编号 I.7.3
层
\[
Rf_{K!}S_W\in D\Bigl([*/G_b(E)]\times\prod_{i\in I}\Spd\dot E_i,\Lambda\Bigr)
\]
带有部分 Frobenii，从而下降为
\[
D\Bigl([*/G_b(E)]\times\prod_{i\in I}\Spd\dot E_i/\varphi_i^\ZZ,\Lambda\Bigr)
\]
的一个对象。
\end{corollary}

\srcpage{32}该对象落在满子范畴
\[
D(G_b(E),\Lambda)^{B\prod_{i\in I}W_{E_i}}\subset D\Bigl([*/G_b(E)]\times\prod_{i\in I}\Spd\dot E_i/\varphi_i^\ZZ,\Lambda\Bigr)
\]
之中，并且它到 $D(G_b(E),\Lambda)$ 的限制是紧的。特别地，对 $G_b(E)$ 的任意容许表示 $\rho$，对象
\[
R\Hom_{G_b(E)}(Rf_{K!}S_W,\rho)\in D(\Lambda)^{B\prod_{i\in I}W_{E_i}}
\]
是 $\prod_{i\in I}W_{E_i}$ 在 $\Lambda$-模的完美复形上的一个表示。对 $K$ 取余极限，就得到一个容许 $G(E)$-表示的复形
\[
\varinjlim_K R\Hom_{G_b(E)}(Rf_{K!}S_W,\rho),
\]
它带有 $\prod_{i\in I}W_{E_i}$-作用。\yizhu{原文此显示式在 $[*/G_b(E)]$ 的右方括号之前多印了一个右括号（印作 $D([*/G_b(E))]$），使 $D(\,\cdot\,,\Lambda)$ 的括号提前闭合；此为疑似排版笔误，译文按其本意 $D([*/G_b(E)]\times\prod_{i\in I}\Spd\dot E_i/\varphi_i^\ZZ,\Lambda)$ 排印。}{A:p.32 第 3 行 “$D([*/G_b(E))]\times\prod\cdots,\Lambda)$”}



若 $\rho$ 是紧的，则
\[
\varinjlim_K R\Hom_{G_b(E)}(Rf_{K!}S_W,\rho)
\]
作为 $G(E)$-表示的复形也是紧的。

特化到 $I=\{*\}$ 且 $\mu$ 极小，就得到局部 Shimura 簇，而无条件地证明了 [RV14, Proposition 6.1] 中的有限性性质，以及 [RV14, Remark 6.2 (iii)]。我们指出，这些性质似乎无法只用 shtuka 模空间的定义（即不用 $\Bun_G$）得到。

% ============================================================
\section{L-参数叠}
% ============================================================

下面讨论 Langlands 对应的另一侧，即 L-参数（的叠）。这先前已由 Dat–Helm–Kurinczuk–Moss [DHKM20] 与 Zhu [Zhu20] 做过。我们希望定义一个概形，使其 $\Lambda$-点（$\Lambda$ 为 $\ZZ_\ell$-代数）是连续 $1$-上闭链
\[
\varphi:W_E\longrightarrow\hat G(\Lambda).
\]
（这里给 $\hat G$ 配上通常的 $W_E$-作用，它经由 $W_E$ 的一个有限商 $Q$ 分解。如前所讨论，在 $\ZZ_\ell[\sqrt q]$ 上两种作用之间的差别消失了，而此处我们觉得使用标准正规化要方便得多，这样有时就可以利用代数群 $\hat G\rtimes Q$。）

这里似乎有一个不匹配之处：我们要求的是一个代数叠，而连续余圈却是另一回事。有趣的是，连续性条件有一种表述方式能产生概形。即：把 $\Lambda$ 视为“在 $\ZZ_\ell$ 上相对离散”的凝聚 $\ZZ_\ell$-代数。抽象 $\ZZ_\ell$-模 $M$ 通过 $M\mapsto M_{\mathrm{disc}}\otimes_{\ZZ_\ell,\mathrm{disc}}\ZZ_\ell$ 全忠实地嵌入凝聚 $\ZZ_\ell$-模。

\begin{theorem}% 原编号 I.8.1
存在 $\ZZ_\ell$ 上的概形 $\fsZ^1(W_E,\hat G)$，其 $\Lambda$-点（$\Lambda$ 为 $\ZZ_\ell$-代数）是凝聚 $1$-上闭链
\[
\varphi:W_E\longrightarrow\hat G(\Lambda),
\]
其中我们把 $\Lambda$ 视为相对离散的凝聚 $\ZZ_\ell$-代数。概形 $\fsZ^1(W_E,\hat G)$ 是开闭仿射子概形 $\fsZ^1(W_E/P,\hat G)$ 之并，其中 $P$ 取遍 $W_E$ 的\srcpageinl{33}野惯性子群的开子群，并且每个 $\fsZ^1(W_E/P,\hat G)$ 都是 $\ZZ_\ell$ 上的平坦局部完全交，维数为 $\dim G$。
\end{theorem}

这里的要点是：$W_E$ 的惯性子群有一个 $\ZZ_\ell$-因子，而在作此定义时，它可以以有趣的方式映入 $\Lambda$。为证明该定理，我们仿照 [DHKM20] 与 [Zhu20]，通过对驯惯性作离散化来定义离散稠密子群 $W\subset W_E/P$，此时限制映射 $\fsZ^1(W_E/P,\hat G)\to\fsZ^1(W,\hat G)$ 是同构，而后者显然是仿射概形。

我们还可以证明关于 $\hat G$ 在 $\fsZ^1(W_E,\hat G)$ 上作用的进一步结果，更确切地说，是关于每个 $\fsZ^1(W_E/P,\hat G)$ 的。对这个结果，我们需要对 $\ell$ 作一个非常小的假设。

\begin{theorem}% 原编号 I.8.2
假设 $\ell$ 不整除 $\pi_0Z(G)$ 的阶（等价地，$\ell$ 不整除 $\pi_1(\hat G)_{\mathrm{tor}}$ 的阶）。则对 $i>0$ 有 $H^i(\hat G,\cO(\fsZ^1(W_E/P,\hat G)))=0$，并且不变式 $\cO(\fsZ^1(W_E/P,\hat G))^{\hat G}$ 的构造与任意基变换交换。代数 $\cO(\fsZ^1(W_E/P,\hat G))^{\hat G}$ 允许用巡游算子给出的显式表示，
\[
\cO\bigl(\fsZ^1(W_E/P,\hat G)\bigr)^{\hat G}=\varinjlim_{(n,F_n\to W)}\cO\bigl(\fsZ^1(F_n,\hat G)\bigr)^{\hat G},
\]
其中余极限取遍自由群 $F_n$ 到 $W\subset W_E/P$ 的所有映射，而 $\fsZ^1(F_n,\hat G)\cong\hat G^n$ 带有同时扭共轭的 $\hat G$-作用。

此外，$\infty$-范畴 $\fsPerf(\fsZ^1(W_E/P,\hat G)/\hat G)$ 在锥与收缩之下由 $\Rep(\hat G)\to\fsPerf(\fsZ^1(W_E/P,\hat G)/\hat G)$ 的像生成，而 $\operatorname{Ind}\fsPerf(\fsZ^1(W_E/P,\hat G))$ 等价于 $\operatorname{Ind}\fsPerf(*/\hat G)$ 上 $\cO(\fsZ^1(W_E/P,\hat G))$-模范畴。
\end{theorem}

上述所有结果在 $\QQ_\ell$-系数下也成立，且无需对 $\ell$ 作假设。

在 $\QQ_\ell$-系数下这些结果是简单的，因为 $\hat G$ 的表示论是半单的。然而在 $\ZZ_\ell$-系数下，这些结果相当微妙，我们需要深入约化群的模表示论。事实上，对“沿约化群的嵌入 $H\subset G$ 作限制保持具有好滤过的表示”这一现象，我们给出了新的视角（以及一些新的例子）。

\begin{theorem}% 原编号 I.8.3
设 $G$ 为代数闭域 $L$（特征 $\ell$）上的约化群。设 $P$ 为阶与 $\ell$ 互素的有限可解群，作用在 $G$ 上。不动点群 $H=G^P$ 是光滑线性代数群，$H^\circ$ 约化，且 $\pi_0H$ 的阶与 $\ell$ 互素。

在此情形下，对 $G$ 的任意具有好 $G$-滤过的表示 $V$，其限制 $V|_{H^\circ}$ 也具有好 $H^\circ$-滤过。
\end{theorem}

Levi 子群的情形是经典定理，而 $P=\ZZ/2\ZZ$ 的情形曾被称为 Brundan 猜想，并由 [Bru98]、[vdK01] 中的穷举式逐例分析所证明。我们给出一个在一切情形下都一致适用的新证明。

我们还证明：在定理的情形下，$\fsPerf(*/G)\to\fsPerf(*/H)$ 的像在锥与收缩之下生成整个范畴。在最初的版本中，我们通过对所有可能情形的非常显式（且令人疲惫）的分析来证明这一点，但现在已有一个统一的证明。

% ============================================================
\section{L-参数的构造}
% ============================================================

\srcpageinl{34}最后，我们可以讨论 L-参数的构造了。先为简单起见假设 $\Lambda=\ol{\QQ}_\ell$ 并固定 $\sqrt q\in\ol{\QQ}_\ell$，并设 $A\in D(\Bun_G,\ol{\QQ}_\ell)$ 为任意 Schur 不可约对象，即 $\End(A)=\ol{\QQ}_\ell$。例如，$A$ 可以对应于 $G(E)$ 的一个不可约光滑表示，即沿 $[*/G(E)]\hookrightarrow\Bun_G$ 作零延拓。于是，仿照 V. Lafforgue [Laf18]，我们可以如下定义巡游算子。对 $(\hat G\rtimes W_E)^I$ 在 $\ol{\QQ}_\ell$ 上的任意表示 $V$，连同映射
\[
\alpha:\ol{\QQ}_\ell\to V|_{\hat G},\qquad \beta:V|_{\hat G}\to\ol{\QQ}_\ell,
\]
它们限制到对角副本 $\hat G\subset(\hat G\rtimes W_E)^I$ 的作用上，再取元素 $\gamma_i\in W_E$（$i\in I$），就可以定义 $A$ 的自同态
\[
A\xrightarrow{T_\alpha}T_V(A)\xrightarrow{(\gamma_i)_{i\in I}}T_V(A)\xrightarrow{T_\beta}A,
\]
从而定义了 $\ol{\QQ}_\ell$ 中的一个元素。形式体系既已齐备，下述结果本质上归功于 V. Lafforgue [Laf18, Proposition 11.7]。

\begin{proposition}% 原编号 I.9.1
存在唯一的连续半单 L-参数
\[
\varphi_A:W_E\longrightarrow\hat G(\ol{\QQ}_\ell),
\]
使得对如上所有的 $(I,V,\alpha,\beta,(\gamma_i)_{i\in I})$，巡游算子
\[
A\xrightarrow{T_\alpha}T_V(A)\xrightarrow{(\gamma_i)_{i\in I}}T_V(A)\xrightarrow{T_\beta}A
\]
等于乘上标量
\[
\ol{\QQ}_\ell\xrightarrow{\alpha}V\xrightarrow{(\varphi_A(\gamma_i))_{i\in I}}V\xrightarrow{\beta}\ol{\QQ}_\ell.
\]
\end{proposition}

注意，事实上巡游算子定义了 $G(E)$ 的 Bernstein 中心中的元素，因为它们定义了恒等函子的自同态。从这个观点出发，我们作如下定义。

\begin{definition}% 原编号 I.9.2
\begin{enumerate}
\item[(i)] $G(E)$ 的 Bernstein 中心是
\[
\mathcal{Z}(G(E),\Lambda)=\pi_0\End\bigl(\mathrm{id}_{D(G(E),\Lambda)}\bigr)=\varprojlim_{K\subset G(E)}\mathcal{Z}\bigl(\Lambda[K\backslash G(E)/K]\bigr),
\]
其中 $K$ 取遍 $G(E)$ 的开 pro-$p$ 子群，而 $\Lambda[K\backslash G(E)/K]=\End_{G(E)}\bigl(\operatorname{c-Ind}_K^{G(E)}\Lambda\bigr)$ 是级 $K$ 的 Hecke 代数。
\item[(ii)] $G$ 的几何 Bernstein 中心是
\[
\mathcal{Z}^{\mathrm{geom}}(G,\Lambda)=\pi_0\End\bigl(\mathrm{id}_{\fslis(\Bun_G,\Lambda)}\bigr).
\]
在 $\mathcal{Z}^{\mathrm{geom}}(G,\Lambda)$ 内部，令 $\mathcal{Z}^{\mathrm{geom}}_{\mathrm{Hecke}}(G,\Lambda)$ 为所有与 Hecke 算子交换的自同态 $f:\mathrm{id}\to\mathrm{id}$ 构成的子环，即对一切 $V\in\Rep(\hat G^I)$ 与 $A\in\fslis(\Bun_G,\Lambda)$，都有 $T_V(f(A))=f(T_V(A))\in\End(T_V(A))$。
\item[(iii)] \srcpageinl{35}$G$ 的谱 Bernstein 中心是
\[
\mathcal{Z}^{\mathrm{spec}}(G,\Lambda)=\cO\bigl(\fsZ^1(W_E,\hat G)_\Lambda\bigr)^{\hat G},
\]
即商叠 $\fsZ^1(W_E,\hat G)_\Lambda/\hat G$ 上整体函数之环。
\end{enumerate}
\end{definition}


包含 $D(G(E),\Lambda)\hookrightarrow\fslis(\Bun_G,\Lambda)$ 诱导出代数映射 $\mathcal{Z}^{\mathrm{geom}}(G,\Lambda)\to\mathcal{Z}(G(E),\Lambda)$。

现在，巡游算子的构造连同定理 I.8.2 蕴含下述结果。这里 $\Lambda$ 是 $\ZZ_\ell[\sqrt q]$-代数，使得 $\pi_0Z(G)$ 的阶在 $\Lambda$ 中可逆。

\begin{proposition}% 原编号 I.9.3
存在典范映射
\[
\mathcal{Z}^{\mathrm{spec}}(G,\Lambda)\longrightarrow\mathcal{Z}^{\mathrm{geom}}_{\mathrm{Hecke}}(G,\Lambda)\subset\mathcal{Z}^{\mathrm{geom}}(G,\Lambda),
\]
特别地，存在一个映射
\[
\Psi_G:\mathcal{Z}^{\mathrm{spec}}(G,\Lambda)\longrightarrow\mathcal{Z}(G(E),\Lambda).
\]
\end{proposition}

上面关于 L-参数的构造于是便是这个 Bernstein 中心上的映射的推论。这样的整映射的存在性，在 $G=\mathrm{GL}_n$ 的情形归功于 Helm–Moss [HM18]。

\begin{remark}% 原编号 I.9.4
在函数域情形，类似的构造已由 Genestier–Lafforgue [GL17] 给出。Li–Huerta [LH23] 证明了这些构造彼此一致。
\end{remark}

关于对 $\ell$ 的无关性，我们作如下的猜想。为叙述它，注意存在自然的 $\QQ$-代数 $\mathcal{Z}^{\mathrm{spec}}(G,\QQ)$，它对任意 $\ell\neq p$ 基变换到 $\QQ_\ell$ 都给出 $\mathcal{Z}^{\mathrm{spec}}(G,\QQ_\ell)$；事实上，可以取 L-参数叠上对离散拓扑连续（即在 $W_E$ 的某个开子群上平凡）的整体函数；另可参看 [DHKM20]。

\begin{conjecture}% 原编号 I.9.5
存在（必然唯一的）映射 $\mathcal{Z}^{\mathrm{spec}}(G,\QQ(\sqrt q))\to\mathcal{Z}(G(E),\QQ(\sqrt q))$，使得对任意 $\ell\neq p$，基变换到 $\QQ_\ell$ 后它复原复合映射
\[
\mathcal{Z}^{\mathrm{spec}}(G,\QQ_\ell(\sqrt q))\longrightarrow\mathcal{Z}^{\mathrm{geom}}(G,\QQ_\ell(\sqrt q))\longrightarrow\mathcal{Z}(G(E),\QQ_\ell(\sqrt q)).
\]
\end{conjecture}

这将保证我们所构造的 L-参数在相关意义下与 $\ell$ 无关。关于这个映射及其与稳定 Bernstein 中心的关系，Haines [Hai14] 提出了进一步的猜想（另见 [BKV15]、[SS13, Section 6]）。特别地，猜测当 $G$ 拟分裂时，映射 $\Psi_G$ 是单射，且其像可刻画为 $G(E)$ 的 Bernstein 中心中那些对应的分布在该稳定共轭下不变的元素。

也可以借局部 shtuka 的模空间来构造到 Bernstein 中心的映射，如下所述。为简单起见，我们仍只对 $\Lambda$ 为某个 $\ZZ/\ell^n$-代数的情形讨论。给定如上 $I$ 与 $V$，我们可以考虑前面考虑过的空间 $\Sht_{(G,b,\mu_\bullet),K}$ 的一个变体 $\Sht_{(G,1,V),K}$：其中界由 $V$ 的支撑给出，并固定元素 $b=1$。它们带有 étale 周期映射
\[
\pi_K:\Sht_{(G,1,V),K}\longrightarrow\Gr^{\mathrm{tw}}_{G,\prod_{i=1}^{n}\Spd\dot E}
\]
与一个反常层 $S_V$。当限制到几何对角线
\[
\bar x:\Spd\dot E\longrightarrow\prod_{i=1}^{n}\Spd\dot E
\]
上时，\srcpageinl{36}它们成为相应的带一条leg的 shtuka 模空间
\[
f^\Delta_K:\Sht_{(G,1,V|_{\hat G})}\longrightarrow\Spd\dot E,
\]
并带有层 $S_{V|_{\hat G}}$。层 $S_{V|_{\hat G}}$ 带有映射 $\alpha$（相应地 $\beta$），其定义域（相应地靶）是层 $i_*\Lambda$，其中 $i:G(E)/K=\Sht_{(G,1,\ol{\QQ}_\ell),K}\hookrightarrow\Sht_{(G,1,V|_{\hat G}),K}$ 是无leg shtuka 的子空间。这就产生了一个自同态
\begin{align*}
\operatorname{c-Ind}_K^{G(E)}\Lambda &\xrightarrow{\alpha}Rf^\Delta_{K!}S_{V|_{\hat G}}=(Rf_{K!}S_V)_{\bar x}\\
&\xrightarrow{(\gamma_i)_{i\in I}}(Rf_{K!}S_V)_{\bar x}=Rf^\Delta_{K!}S_{V|_{\hat G}}\xrightarrow{\beta}\operatorname{c-Ind}_K^{G(E)}\Lambda.
\end{align*}
这里 $(\gamma_i)_{i\in I}$ 的作用由推论 I.7.3 定义。由定义可知，把先前的构造施用于表示 $\operatorname{c-Ind}_K^{G(E)}\Lambda$，得到的正是这个自同态。注意这些自同态是 $G(E)$-等变的，因此定义了 Hecke 代数
\[
\Lambda[K\backslash G(E)/K]=\End_{G(E)}\bigl(\operatorname{c-Ind}_K^{G(E)}\Lambda\bigr)
\]
中的元素；事实上这些元素是中心的（与先前的构造比较即知）。对 $K$ 取逆向极限，就得到 $G(E)$ 的 Bernstein 中心中的元素。\footnote{当第二作者作他的 Berkeley 讲义 [SW20] 时，这正是我们所设想的巡游算子的构造。注意这里的关键一步是：局部 shtuka 模空间的上同调定义了 $(\Div^1)^I$ 上的一个局部系统。然而，如何纯粹用 shtuka 模空间来证明这一点并不清楚。在整体函数域情形，这一结果已由 Xue [Xue20] 得到。}

关于我们所构造的 L-参数，可以证明下述基本结果。

\begin{theorem}% 原编号 I.9.6
\begin{enumerate}
\item[(i)] 若 $G=T$ 是环面，则 $\pi\mapsto\varphi_\pi$ 就是通常的 Langlands 对应。
\item[(ii)] 对应 $\pi\mapsto\varphi_\pi$ 与扭转相容。
\item[(iii)] 对应 $\pi\mapsto\varphi_\pi$ 与中心特征标相容（参看 [Bor79, 10.1]）。
\item[(iv)] 对应 $\pi\mapsto\varphi_\pi$ 与取反变（contragredient）表示相容（参看 [AV16]）。\yizhuw{原文印作 “passage to congradients”；据所引 [AV16]（Adams–Vogan，\emph{Contragredient representations and characterizing the local Langlands correspondence}）及文意，此处应为 contragredients，系原文拼写笔误。译文按文意译出并保留原拼写。}
\item[(v)] 若 $G'\to G$ 是约化群的映射，诱导出伴随群的同构，$\pi$ 是 $G(E)$ 的不可约光滑表示，而 $\pi'$ 是 $\pi|_{G'(E)}$ 的某个不可约成分，则 $\varphi_{\pi'}$ 是 $\varphi_\pi$ 在诱导映射 $\hat G\to\hat G'$ 下的像。
\item[(vi)] 若 $G=G_1\times G_2$ 是两个群的积，$\pi$ 是 $G(E)$ 的不可约光滑表示，则 $\pi=\pi_1\boxtimes\pi_2$，其中 $\pi_i$ 是 $G_i(E)$ 的不可约光滑表示，且在 $\hat G=\hat G_1\times\hat G_2$ 之下有 $\varphi_\pi=\varphi_{\pi_1}\times\varphi_{\pi_2}$。
\item[(vii)] 若 $G=\Res_{E'|E}G'$ 是某个有限可分扩张 $E'|E$ 上约化群 $G'$ 的 Weil 限制，从而 $G(E)=G'(E')$，则 $G|E$ 的 L-参数与 $G'|E'$ 的 L-参数一致。
\item[(viii)] 对应 $\pi\mapsto\varphi_\pi$ 与抛物诱导相容。
\item[(ix)] 对 $G=\mathrm{GL}_n$ 与超尖点 $\pi$，对应 $\pi\mapsto\varphi_\pi$ 与通常的局部 Langlands 对应 [LRS93]、[HT01]、[Hen00] 一致。
\end{enumerate}
\end{theorem}


注意，把 (viii) 与 (ix) 合起来，就说明对 $\mathrm{GL}_n$ 与一般的 $\pi$，L-参数 $\varphi_\pi$ 就是通常所谓的半单 L-参数。
